\chapter{Introduction}
\label{ch:intro}

Asteroids are leftovers from the formation of the Solar System. Before a spacecraft can land on
one or take a sample, scientists need a detailed photographic map of its surface. Most small
bodies are irregular, non-convex shapes whose ridges and craters hide parts of the surface from
view, and mapping them takes a long time: NEAR Shoemaker, the only spacecraft that ever orbited
433~Eros, spent about a year lowering its orbit step by step to map the body with a single
camera~\cite{dunham2002}.

This project asks how a small fleet of \emph{cooperating} spacecraft should be designed so that
together they photograph as much of an irregular body as possible, without spending more
spacecraft, energy or mission time than needed. Milestone~2 reviewed the literature and
formulated the problem. Milestone~3 implemented \acf{SA} for it, applied it to case studies of
increasing size and difficulty, and validated the results. This milestone (Milestone~4) adds a
\acf{GA} on the same pipeline, validates it in the same way and compares it with \ac{SA}. The
Milestone~3 chapters are kept, updated to the final code, because the \ac{GA} is measured against them.

\section{Objectives}
The objectives of Milestones~3 and~4, following the project brief, are:
\begin{enumerate}
\item Finalise the problem formulation (problem statement, decision variables, constraints and
fitness) and implement it as a modular simulation pipeline (Chapters~\ref{ch:method}
and~\ref{ch:orbital}).
\item Implement \ac{SA} for the problem, including its problem-specific neighbourhood and repair
operators, and visualise the search and the resulting solutions (Chapter~\ref{ch:sa}).
\item Define case studies of varying size and complexity and their input parameters
(Chapter~\ref{ch:cases}).
\item Validate the algorithm on benchmark problems with known optima, against random search,
against results from the literature, and with independent checks of its decisions
(Chapter~\ref{ch:validation}).
\item Report and discuss the results of every case study, and analyse the performance of the
algorithm and the effect of its parameters (Chapter~\ref{ch:results}).
\item \textbf{Milestone~4:} implement a \ac{GA} with a problem-specific chromosome, selection,
crossover and mutation (Chapter~\ref{ch:ga}), validate it, study its parameters, and compare it
with \ac{SA} on the same case studies and budgets (Chapter~\ref{ch:garesults}).
\end{enumerate}

\section{What changed since Milestone 2}
\begin{itemize}
\item Only the orbital layer is optimised. Orbits are circular (eccentricity and argument of
periapsis removed) and every spacecraft carries the same narrow camera (scanner type removed).
\item The mission duration $T$ is a decision variable that is minimised, instead of a fixed three
days, and the number of spacecraft $N$ is decided by the optimiser through one on/off switch per
spacecraft slot.
\item The model now includes sunlight with cast shadows, eclipses and a battery constraint, and the
objectives are coverage, fleet cost, energy and mission time.
\item The method was extended to a second, very different body, 216~Kleopatra.
\end{itemize}

\section{Report outline}
Chapter~\ref{ch:background} summarises the literature reviewed in Milestone~2.
Chapter~\ref{ch:method} gives the final problem formulation and Chapter~\ref{ch:orbital} the
simulation that evaluates it. Chapter~\ref{ch:sa} explains the \ac{SA} implementation,
Chapter~\ref{ch:cases} the case studies, Chapter~\ref{ch:validation} the validation and
Chapter~\ref{ch:results} the results and the analysis of the algorithm.
Chapter~\ref{ch:ga} explains the \ac{GA} implementation and Chapter~\ref{ch:garesults} its
validation, parameter study and comparison with \ac{SA}.
Chapters~\ref{ch:conclusion} and~\ref{ch:future} conclude and outline the next milestones.
Appendix~\ref{app:code} explains how to run the code.
